In this chapter I will disclose the general aspects of the approach regarding the insertion of the Edge Tracing 
instructions inside an already compiled binary.\\
There are three main steps to be fullfilled to have a binary that is able to trace Edge Coverage and 
    to be compatible with AFL/AFL++:
    \begin{enumerate}
        \item Insertion of the forkserver invocation prior the main (\cref{sec:forkserver_insertion}).
        \item Search of instrumentation points to trace Edge Coverage.
        \item Intrument the previously found points with Edge Coverage tracing instructions.
    \end{enumerate}

\section{Forkserver Insertion}\label{sec:forkserver_insertion}
    This has been one of the most crucial points of this work.
    First of all I have tried to instrument the forkserver function directly inline.
    This approach present many problems. The biggest difficulty is the fact that a forkserver needs some library functions 
    to work properly :
    \begin{itemize}
        \item \texttt{read/write} to implement the communication protocol with the fuzzer.
        \item \texttt{fork/waitpid} to implement the forkserver main logic.
    \end{itemize}
    These functions could be part of the target binary (if the binary already uses these functions) or not.
    In the best case the binary already uses these functions and they are available. 
    In this case a call to these functions could be instrumented at any time taking into consideration the type of 
    linking with which the binary has been comppiled.
    \begin{itemize}
        \item In case of \textit{Static Linking} the function instructions are already inside the target application, 
            because during the compilation the linker inserted it. So a call to these functions could be inserted finding 
            the function address and then instrument a call instruction to that address.
        \item In case of \textit{Dynamic Linking} the function instructions are not inside the target application when 
            the application is launched, but the function addresses are dynamically resolved by the dynamic linker at need. 
            In this case two new section must be inserted (\texttt{.plt} and \texttt{.got} sections).
            As can be seen in \cref{fig:dynamic_linking_example} 
    \end{itemize}
    \begin{figure}[ht]
        \centering
        \begin{minipage}[h]{0.48\textwidth}
            \centering
            \caption{An example of printf function resolution using Dynamic Linking and Immediate Binding}
            \includegraphics[width=150pt]{dynamicFunctionResolution.png}
            \label{fig:dynamic_linking_example}
        \end{minipage}
        \hfill
        \begin{minipage}[h]{0.48\textwidth}
            \begin{small}
            In \cref{fig:dynamic_linking_example} can be seen an example of dynamic resolution of \texttt{printf} function.
            In case of Dynamic Linking the printf function call in the main is not done to the actual printf libc function.
            Instead there is a specific function in the .plt section usually called  \texttt{printf@plt}.
            Into the .plt section each library function is replaced by the following lines of code
            \begin{lstlisting}
jmp    QWORD PTR [rip + offset_to_printf@GOT]
push   0x0
jmp    <plt_starting_address>
            \end{lstlisting}
            Into the .got section there are the resolved addresses of the external functions.
            If the .got actually contains the resolved address of the printf (the library address), the first intruction 
            will jump into the library.
            The second and third instructions are used in case of \textbf{Lazy Binding}. In this case the GOT entry 
            does not contain the address of the library function when the main start, but contain the next address of the plt entry.
            At the \texttt{plt\_starting\_address} there is the plt resolver that invoke the dynamic linker that search for the 
            address of printf function into libc.so library. 
            \end{small} 
        \end{minipage}
    \end{figure}
    If the functions required by the forkserver are part of the target binary, their calls can be easily instrumented. 
    Depending on whether the binary is statically or dynamically linked, this can be done either by finding the address 
    of the function and calling it directly, or by finding the address of its .plt entry and calling it, 
    leaving the dynamic linker to resolve the function. 
    The problem arises when the original binary does not include these functions. 
    In this case, it is almost impossible to modify the .plt and .got sections to include the new functions 
    while preserving the binary layout and correctness in the general case. In fact, in some exceptional cases, 
    padding may provide enough space in the .got section for a new entry, but this cannot be assumed in the general case. 
    Therefore, this first approach has been completely discarded.

    Since the forkserver is more complicated than the simple shared map update due to the presence of many controls on 
    environment variables and communication with the fuzzer, 
    the best approach is to inject the compiled forkserver code inside the target.
    This has been done through a \textit{Shared Onject Injection}.
    In order to achieve a Shared Object Injection the \textit{.dynamic} section of the target binary is modified.
    In particular :
    \begin{itemize}
        \item A \textit{DT\_NEEDED} entry is added with the name of the shared object to be injected. \footnote{ In reality 
            the shared object name is a string that is written into the \textit{.dynstr} section and into the 
            \textit{.dynamic} section is specified the offset of the string written into \textit{.dynstr}}
        \item A \textit{RUNPATH} entry is added containing the path in which is saved the actual shared object.
    \end{itemize}

    The shared object injected into the target binary, contains the code of the forkserver that must be executed 
    prior the main of the binary.
    In order to achieve the fact that a function of a shared object injected into a binary is executed before the main 
    of the binary, that function must be declared as ``constructor''.
    This means prepend to the function the string \texttt{\_\_attribute\_\_((constructor))}.

